\documentclass[10pt,a4paper]{amsart}
\usepackage[margin=19mm]{geometry}
\usepackage{amsmath,amssymb,mathtools}
\usepackage[scr=rsfs,cal=euler]{mathalfa}
\usepackage{xfrac}
\usepackage[unicode,hidelinks,bookmarks=false]{hyperref}
\newcommand{\ie}{i.e.\ }
\newcommand{\cf}{cf.\ }
\newcommand{\resp}{respectively\ }
\newcommand{\Subsection}{\subsection}
\providecommand{\nobreakdash}{\nobreak\hbox{-}}
\setlength{\emergencystretch}{2em}
\multlinegap0pt
\usepackage{bm}
\usepackage{bbm}
\usepackage{dsfont}
\usepackage{xspace}
\usepackage{cancel}
\usepackage{mathtools}
\usepackage{amsthm}
\makeatletter
\@ifundefined{theorem}{\newtheorem{theorem}{Theorem}}{}
\@ifundefined{claim}{\newtheorem{claim}[theorem]{Claim}}{}
\makeatother
\title[Orbit pseudometrics and a universality property of the Gromo]{Orbit pseudometrics and a universality property of the Gromov-Hausdorff distance}
\author{Ond\v{r}ej Kurka}
\date{Source: arXiv:2204.08375v2 (CC BY 4.0)}
\begin{document}
\maketitle

\noindent\emph{Source and licence.} Orbit pseudometrics and a universality property of the Gromov-Hausdorff distance. Version arXiv:2204.08375v2 (CC BY 4.0), distributed under the Creative Commons Attribution 4.0 International licence, as recorded in the arXiv licence metadata for this exact version and shown on the abstract page https://arxiv.org/abs/2204.08375v2. Full rights notice: licenses/arxiv-2204.08375-CC-BY-4.0.txt.\par\smallskip

\noindent\textit{Editorial scope.} Counted unit: the Claim whose source label is \texttt{cl6} in the article (source0.tex lines 330--333, source label \texttt{cl6}), delivered together with the article's own complete original proof of it (source0.tex lines 335--369, one \texttt{proof} environment of 35 lines). No mathematics is written, repaired or re-derived: statement and proof are the article's own text line for line. Required context retained uncounted: cl5 (lines 299--307). Every definition, display and label of those lines is retained unchanged. Omitted from this package and not redistributed: 1--298, 308--329, 334--334, 370--662. Nothing inside a retained excerpt is omitted or altered: every retained body is delivered exactly as the article's lines are written, so no omission receipt of any kind accompanies this package. Citations: the retained excerpts carry no citation command at all, so no bibliography entry of the article is transcribed and no reference is redistributed. Reference adaptation: none was needed and none was made; every reference inside the delivered unit resolves inside it, so no unresolved reference or citation is produced. Printed slips kept literal: no printed word, symbol, hypothesis or number of the counted statement or of its proof was altered. Layout and class: the article's own class is \texttt{amsart} with packages including mathtools; the wrapper is the lane's pinned \texttt{amsart} editorial wrapper at 10pt on A4, which restates the theorem-like environments the retained text needs and adds the article's own macro definitions that the retained text needs (none), with extra wrapper packages (bm, bbm, dsfont, xspace, cancel, mathtools). The delivered numbering is the unit's own. Assets: no retained span contains a figure environment, a graphics-inclusion command, a PDF-page inclusion, a table or a TikZ picture; no image file is copied into this package, the package declares no asset dependency and no image file is redistributed with it. Licence: version arXiv:2204.08375v2 of this article is distributed under the Creative Commons Attribution 4.0 International licence, as recorded in the arXiv licence metadata for this exact version and shown on the abstract page https://arxiv.org/abs/2204.08375v2. A full rights notice is delivered at licenses/arxiv-2204.08375-CC-BY-4.0.txt. Characters: every retained excerpt is pure ASCII, so the delivered engine, XeLaTeX with its default Latin Modern text font, reports no missing character.
\begin{claim} \label{cl5}
Let $ p \in X, w \in W_{p} $, and let $ A(w, p) = \{ m(w', w) : w' \in W_{p} \} $. Let further $ B \subseteq [0, \infty) $ be such that $ \varrho_{H}(B, A(w, p)) \leq 1 $ and $ \beta_{l} $ be defined by $ \beta_{l} = \inf (B \cap (80 \cdot 2^{l}, \infty)) $ for $ l \in \mathbb{N} $.

If $ w \in Z \times \{ k \} $, then $ \liminf_{l \to \infty} (100 \cdot 2^{l} - \beta_{l}) \in [100 \cdot 2^{k} - 1, 100 \cdot 2^{k} + 1] $.

If $ w \in Z \times [0, 1] \times \{ k \} $, then $ \liminf_{l \to \infty} (100 \cdot 2^{l} - \beta_{l}) \in [90 \cdot 2^{k} - 2, 90 \cdot 2^{k} + 1] $.

Consequently, if $ \varrho_{H}(A(w, p), A(w', q)) \leq 1 $ for some $ p, q \in X, w \in W_{p}, w' \in W_{q} $, then $ w \in Z \times \{ k \} $ is equivalent to $ w' \in Z \times \{ k \} $ and $ w \in Z \times [0, 1] \times \{ k \} $ is equivalent to $ w' \in Z \times [0, 1] \times \{ k \} $.
\end{claim}

\begin{claim} \label{cl6}
For all $ p, q \in X $,
$$ \varrho_{G, d}(p, q) \leq 2 \varrho_{GH}(Y_{p}, Y_{q}). $$
\end{claim}

\begin{proof}
Given $ r > \varrho_{GH}(Y_{p}, Y_{q}) $, we need to show that $ 2r \geq \varrho_{G, d}(p, q) $. Since $ r > \varrho_{GH}(Y_{p}, Y_{q}) = \varrho_{GH} (W_{p}, W_{q}) $, there is a correspondence $ \mathcal{R} $ between $ W_{p} $ and $ W_{q} $ such that
$$ a \mathcal{R} a' \; \& \; b \mathcal{R} b' \quad \Rightarrow \quad |m(a, b) - m(a', b')| \leq 2r. $$
As $ \varrho_{G, d}(p, q) \leq 1 $, we can assume that $ 2r < 1 $. In such a case, due to Claim~\ref{cl5}, elements of $ Z \times \{ k \} $ may correspond only to elements of $ Z \times \{ k \} $ itself, and the same holds for $ Z \times [0, 1] \times \{ k \} $. Consequently, the relations
$$ \mathcal{R}_{k} = \big\{ (z, z') \in Z \times Z : (z, k) \mathcal{R} (z', k) \big\}, \quad k \in \mathbb{N}, $$
are correspondences on $ Z $.

It is easy to check that
$$ z \mathcal{R}_{k_{1}} z_{1} \; \& \; z \mathcal{R}_{k_{2}} z_{2} \quad \Rightarrow \quad \delta_{Z}(z_{1}, z_{2}) \leq 2^{-\min \{ k_{1}, k_{2} \} } \cdot 2r, $$
as $ (z, k_{1}) \mathcal{R} (z_{1}, k_{1}), (z, k_{2}) \mathcal{R} (z_{2}, k_{2}) $ implies $ 2r \geq |m((z, k_{1}), (z, k_{2})) - m((z_{1}, k_{1}), (z_{2}, k_{2}))| = | 100 \cdot |2^{k_{1}} - 2^{k_{2}}| + 2^{\min \{ k_{1}, k_{2} \} } \delta_{Z}(z, z) - 100 \cdot |2^{k_{1}} - 2^{k_{2}}| - 2^{\min \{ k_{1}, k_{2} \} } \delta_{Z}(z_{1}, z_{2}) | = 2^{\min \{ k_{1}, k_{2} \} } \delta_{Z}(z_{1}, z_{2}) $.

It follows that for every $ z \in Z $, there is $ I(z) \in Z $ such that
$$ z \mathcal{R}_{k} z' \quad \Rightarrow \quad \delta_{Z}(I(z), z') \leq 2^{-k} \cdot 2r $$
(if we consider $ S_{j} = \bigcup_{k=j}^{\infty} z \mathcal{R}_{k} $, then the diameter of $ S_{j} $ is at most $ 2^{-j} \cdot 2r $, and we can take $ \{ I(z) \} = \bigcap_{j=1}^{\infty} \overline{S_{j}} $). Similarly, for every $ z \in Z $, there is $ J(z) \in Z $ such that
$$ z' \mathcal{R}_{k} z \quad \Rightarrow \quad \delta_{Z}(z', J(z)) \leq 2^{-k} \cdot 2r. $$

We claim that $ I $ and $ J $ are isometries. It is sufficient to consider $ I $ only. Let $ z_{1}, z_{2} \in Z $, and let $ k \in \mathbb{N} $ be arbitrary. If we choose $ z'_{1}, z'_{2} $ such that $ z_{1} \mathcal{R}_{k} z'_{1}, z_{2} \mathcal{R}_{k} z'_{2} $, then $ \delta_{Z}(I(z_{1}), z'_{1}) \leq 2^{-k} \cdot 2r, \delta_{Z}(I(z_{2}), z'_{2}) \leq 2^{-k} \cdot 2r $. At the same time, from $ (z_{1}, k) \mathcal{R} (z'_{1}, k), (z_{2}, k) \mathcal{R} (z'_{2}, k) $ we get $ | m((z_{1}, k), (z_{2}, k)) - m((z'_{1}, k), (z'_{2}, k)) | \leq 2r $, i.e., $ 2^{k} \cdot | \delta_{Z}(z_{1}, z_{2}) - \delta_{Z}(z'_{1}, z'_{2}) | \leq 2r $. By the triangle inequality, $ | \delta_{Z}(z_{1}, z_{2}) - \delta_{Z}(I(z_{1}), I(z_{2})) | \leq 3 \cdot 2^{-k} \cdot 2r $. Since $ k \in \mathbb{N} $ was arbitrary, we finally see that $ \delta_{Z}(I(z_{1}), I(z_{2})) = \delta_{Z}(z_{1}, z_{2}) $.

In order to show that $ I $ is surjective, we further claim that $ J $ is inverse to $ I $. Let $ z \in Z $, and let $ k \in \mathbb{N} $ be arbitrary. If we choose $ z' $ such that $ z \mathcal{R}_{k} z' $, then $ \delta_{Z}(I(z), z') \leq 2^{-k} \cdot 2r $ and $ \delta_{Z}(z, J(z')) \leq 2^{-k} \cdot 2r $. As $ \delta_{Z}(J(I(z)), J(z')) = \delta_{Z}(I(z), z') $, the triangle inequality provides $ \delta_{Z}(J(I(z)), z) \leq 2 \cdot 2^{-k} \cdot 2r $. Since $ k \in \mathbb{N} $ was arbitrary, we see that $ \delta_{Z}(J(I(z)), z) = 0 $.

So, we have shown that $ I $ is a surjective isometry on $ Z $, and it follows that $ I $ is the extension of $ I_{h} $ for some $ h \in G $.

Now, let $ k \in \mathbb{N} $ be arbitrary. If we choose $ z $ such that $ (1_{G}, p) \mathcal{R}_{k} z $ and $ g, x, u $ such that $ (1_{G}, p, 0, k) \mathcal{R} (g, x, u, k) $, then
$$ \big| m((1_{G}, p, k), (1_{G}, p, 0, k)) - m((z, k), (g, x, u, k)) \big| \leq 2r. $$
The involved distances are $ 10 \cdot 2^{k} $ and $ u + 10 \cdot 2^{k} + 2^{k} \delta_{Z}(z, (g, x)) $, hence we get
$$ u + 2^{k} \delta_{Z}(z, (g, x)) \leq 2r. $$
As $ (g, x, u, k) \in W_{q} $, we have $ \varphi_{k}(gq, x) \leq u \leq 1 $, and thus
$$ \varphi_{k}(gq, x) \leq 2r. $$
Since $ \delta_{Z}((h, p), z) = \delta_{Z}(I(1_{G}, p), z) \leq 2^{-k} \cdot 2r $, the triangle inequality provides $ \delta_{Z}((h, p), (g, x)) \leq 2^{-k} \cdot 2r + 2^{-k} \cdot 2r $, i.e.,
$$ \gamma(h, g) \leq 2^{-k} \cdot 4r, \quad \delta_{X}(p, x) \leq 2^{-k} \cdot 4r. $$

So, we have seen that for every $ k \in \mathbb{N} $, there are $ g_{k} \in G $ and $ x_{k} \in X $ such that
$$ \varphi_{k}(g_{k}q, x_{k}) \leq 2r, \quad \gamma(h, g_{k}) \leq 2^{-k} \cdot 4r, \quad \delta_{X}(p, x_{k}) \leq 2^{-k} \cdot 4r. $$
We have $ g_{k} \to h $, $ x_{k} \to p $, and for every $ \varphi \in \Phi' $, we get $ \varphi(hq, p) \leq 2r $, as it is lower semicontinuous and $ \varphi = \varphi_{k} $ infinitely many times. It follows that $ \varrho_{G, d}(q, p) \leq d(hq, p) \leq 2r $.
\end{proof}

\end{document}
